\documentclass[11pt]{article}
\usepackage[a4paper,margin=27mm]{geometry}
\usepackage{amsmath,amssymb,amsthm,mathtools,xcolor,booktabs,array}
\newcommand{\PROVED}{\textcolor{green!40!black}{\textbf{PROVED}}}
\newcommand{\OPEN}{\textcolor{red!70!black}{\textbf{OPEN}}}
\newcommand{\AUDIT}{\textcolor{orange!70!black}{\textbf{AUDIT}}}
\newtheorem{theorem}{Theorem}
\newtheorem{lemma}{Lemma}
\newtheorem{proposition}{Proposition}
\title{Uniform observability restricted to a URGT$_1$ minimizer (v129)}
\date{14 August 2026}

\begin{document}
\maketitle

\section{The generalized Rayleigh problem}

On $E_N=L^2(0,\delta_N)$ write
\begin{align}
 \mathfrak a_N[e]
 &:=\|\mathbf L_Ne\|^2
 =\int_{\mathbb R}W_N(\tau)|\widehat e(\tau)|^2
                    \frac{d\tau}{2\pi},                   \tag{1}\\
 \mathfrak q_N[e]
 &:=\mathcal E_{\Gamma T,N}(e)
 =\int_{\mathbb R}g_\Gamma(\tau)W_N(\tau)|\widehat e(\tau)|^2
                    \frac{d\tau}{2\pi}+\mathfrak t_N[e], \tag{2}
\end{align}
where $\mathfrak t_N\ge0$ is the rank-two Tate form.  The URGT assertion
is $\mathfrak q_N\ge\mathfrak a_N$.  If it fails, put
\begin{equation}
 \lambda_N:=\inf\{\mathfrak q_N[e]:\mathfrak a_N[e]=1\}<1. \tag{3}
\end{equation}

Let
\begin{equation}
 B=\{\tau:g_\Gamma(\tau)<1\}
 \subset\left(-\frac1{\sqrt{12}},\frac1{\sqrt{12}}\right).
                                                                    \tag{4}
\end{equation}

\begin{theorem}[A failure produces an actual minimizing mode --- \PROVED]
If (3) holds, the infimum is attained by a nonzero
$e_N\in\operatorname{Dom}\mathfrak q_N$, which may be normalized by
$\mathfrak a_N[e_N]=1$.  It satisfies the weak Euler--Lagrange equation
\begin{equation}
 \boxed{
 \int_{\mathbb R}(g_\Gamma-\lambda_N)W_N
       \widehat e_N\overline{\widehat h}\frac{d\tau}{2\pi}
 +\mathfrak t_N(e_N,h)=0
 \quad(h\in\operatorname{Dom}\mathfrak q_N).}             \tag{5}
\end{equation}
\end{theorem}

\begin{proof}
For $N=2$ the leakage column is zero, so (3) cannot occur.  Fix $N\ge3$.
Choose one nonzero depth polynomial, for instance depth one, and list its
distinct physical delays as
\[
 u_1<u_2<\cdots<u_m,
 \qquad
 L_{N,1}=\sum_{j=1}^m c_j\tau_{u_j},
 \qquad c_1\ne0.
\]
Project the output onto the interval $u_1+(0,\delta_N)$ and translate this
interval back to $(0,\delta_N)$.  On zero-extended inputs the resulting
operator is
\[
 c_1(I+D_N),
 \qquad
 D_N=\sum_{j=2}^m\frac{c_j}{c_1}S_{u_j-u_1},
\]
where $S_v$ is the causal right shift by $v$.  If
$v_*:=\min_{j\ge2}(u_j-u_1)$ (and $D_N=0$ when $m=1$), then
$D_N^{r}=0$ for $r>\delta_N/v_*$.  Hence
\[
 (I+D_N)^{-1}=\sum_{r=0}^{M_N-1}(-D_N)^r
\]
for a finite $M_N$, and therefore there is a constant $c_N>0$ such that
\begin{equation}
 \mathfrak a_N[e]\ge \|L_{N,1}e\|_2^2
 \ge c_N\|e\|_2^2.                                      \tag{6a}
\end{equation}
Thus the $\mathfrak a_N$-norm and the ordinary $L^2(0,\delta_N)$ norm
are equivalent for each fixed $N$.

The negative part of $g_\Gamma-1$ is bounded and supported in the compact
set $B$.  The map
\begin{equation}
 e\longmapsto
 \mathbf1_B(\tau)W_N(\tau)^{1/2}\widehat e(\tau)           \tag{6}
\end{equation}
from the finite time interval $E_N$ to $L^2(B)$ is compact: its kernel is
square-integrable on $B\times(0,\delta_N)$, and $W_N$ is bounded on $B$.
By (6a) this compactness is also valid in the $\mathfrak a_N$ topology.
Thus the negative part of the closed form
$\mathfrak q_N-\mathfrak a_N$ is relatively compact with respect to its
nonnegative part.  If the form has a negative direction, the min--max
principle produces a negative discrete generalized eigenvalue below the
essential threshold zero.  Its eigenvector attains (3).  Polarization of
$\mathfrak q_N[e_N]=\lambda_N\mathfrak a_N[e_N]$ gives (5).
\end{proof}

The theorem removes irrelevant near-kernel baggage: the only source which
has to be observed is an eigenvector satisfying (5), not an arbitrary
Paley--Wiener function.

\section{Correctly normalized critical resolvent}

The all-depth result v127 is
\begin{equation}
 \left(\frac{\log N}{N}\right)^2W_N(\tau)
 \longrightarrow
 C_{\rm W}(\tau^2+\tfrac14)^{-2}                           \tag{7}
\end{equation}
locally uniformly.  Therefore the critical source resolvent is
\begin{equation}
 \widehat{R_Ne}(\tau)
 :=\frac{N}{\log N}\frac{\widehat e(\tau)}
                         {\tau^2+\frac14}.                  \tag{8}
\end{equation}
The factor is $N/\log N$; its inverse would not match (7).

\begin{theorem}[Uniform minimizer observability --- \OPEN]
There is a function $\omega(R)\downarrow0$, independent of $N$, such that
every normalized minimizer from Theorem 1 with $\lambda_N<1$ satisfies
\begin{equation}
 \boxed{
 \int_{|\tau|>R}|\widehat{R_Ne_N}(\tau)|^2\frac{d\tau}{2\pi}
 \le\omega(R)\mathfrak a_N[e_N]=\omega(R).}               \tag{9}
\end{equation}
Equivalently,
\begin{equation}
 \int_{|\tau|>R}
 \left(\frac{N}{\log N}\right)^2
 \frac{|\widehat e_N(\tau)|^2}{(\tau^2+\frac14)^2}
 \frac{d\tau}{2\pi}\le\omega(R).                        \tag{10}
\end{equation}
\end{theorem}

Assuming (9), the resolvent compactness lemma of v128 applies.  The only
possible limits are
\begin{equation}
 \frac{c_0+c_1i\tau}{\tau^2+\frac14}.                      \tag{11}
\end{equation}
The constant mode has Gamma ratio $>1+1/45$ and the linear mode has ratio
$3$, contradicting $\lambda_N<1$.  Hence (9) rules out every unbounded
sequence of minimizing infractors.

\section{What (5) does and does not give automatically}

Splitting (5) at the compact band $B$ gives
\begin{align}
 &\int_{B^c}(g_\Gamma-\lambda_N)W_N
       \widehat e_N\overline{\widehat h}\frac{d\tau}{2\pi}
 +\mathfrak t_N(e_N,h)                                    \notag\\
 &\qquad=
 \int_B(\lambda_N-g_\Gamma)W_N
       \widehat e_N\overline{\widehat h}\frac{d\tau}{2\pi}.  \tag{12}
\end{align}
The right side is a compact, band-limited functional.  Thus the minimizer
is generated by compact-frequency data through the inverse of the
positive left side.  Observability (9) would follow from the uniform
restricted inverse estimate
\begin{equation}
 \boxed{
 \|R_Nh\|_2
 \le C\left|
  \mathbf1_BW_N^{1/2}\widehat h
 \right|_2
 \quad\text{on the generalized eigenspace of (5)},}       \tag{13}
\end{equation}
with $C$ independent of $N$ and $\lambda_N<1$.

Equation (13) is not a consequence of (5) alone: the multiplier $W_N$ can
have real zeros (already for $N=3$ the depth-one polynomial is a cosine).
A proof must use time limitation of $h$, the complete all-depth column,
and the fact that $h$ is an eigenvector.  Assuming that no infractor exists
would make (13) circular.

\begin{center}
\begin{tabular}{>{\raggedright\arraybackslash}p{.58\linewidth}
                >{\centering\arraybackslash}p{.14\linewidth}
                >{\raggedright\arraybackslash}p{.18\linewidth}}
\toprule
Statement & Status & Location \\
\midrule
Compactness of the negative frequency part & \PROVED & (4)--(6) \\
Existence of a genuine minimizing infractor & \PROVED & Theorem 1 \\
Exact Euler--Lagrange equation & \PROVED & (5) \\
Correct critical-resolvent normalization & \PROVED & (7)--(8) \\
Uniform observability on minimizing modes & \OPEN & (9)/(13) \\
URGT$_1$, $G$, row (d) & \OPEN & after observability and finite audit \\
\bottomrule
\end{tabular}
\end{center}

\end{document}
